\documentclass[12pt]{article}
\usepackage[margin=2.5cm]{geometry}
\usepackage{parskip}
\usepackage{booktabs}
\usepackage{amsmath}
\usepackage{microtype}
\usepackage{xcolor}
\usepackage[colorlinks=true,linkcolor=blue,citecolor=blue]{hyperref}

\newenvironment{revcomment}{%
  \begin{quote}\itshape
}{%
  \end{quote}
}

\begin{document}

\begin{center}
  {\Large\bfseries Response to Reviewers}\\[1ex]
  {\normalsize Manuscript: ``Bidirectional Meeting Point Assignment and Dynamic Routing\\
  for Many-to-Many Demand-Responsive Transit with Simulated Passenger Response''}\\[0.5ex]
  {\normalsize Journal: \textit{Transportation Research Part A: Policy and Practice}}\\[0.5ex]
  {\normalsize Review date: 2026-04-28 \quad Revision date: 2026-04-28}
\end{center}

\bigskip

We thank the reviewer for the thorough and technically rigorous review. The review identified important gaps between the paper's claims and its actual methodology, as well as inconsistencies in experimental reporting. We have substantially revised the manuscript to close these gaps. Below we address each concern point by point.

\bigskip\hrule\bigskip

\section*{CRITICAL 1: Core ``co-optimization with passenger choice'' claim vs.\ actual algorithm}

\begin{revcomment}
The paper claims ``co-optimization of meeting point assignment, routing, and passenger choice,'' but Sections 4/6 state that bundle selection minimizes only deterministic incremental routing cost, with $P_{accept}$ used only for post-offer Bernoulli sampling. This is ``route-first + ex-post stochastic acceptance simulation,'' not passenger-choice-aware co-optimization. Also, outside option $U_0=0$ with only negative utility coefficients means $P_{accept} \leq 0.5$, yet DoorToDoor acceptance is 0.610.
\end{revcomment}

\textbf{Response.} We accept this criticism fully. We have revised the manuscript to accurately describe what the framework does, rather than over-claiming. Specific changes:

\begin{enumerate}
  \item \textbf{Title:} Changed from ``Co-optimization of Bidirectional Meeting Point Assignment and Dynamic Routing for Many-to-Many Demand-Responsive Transit with Passenger Choice'' to ``Bidirectional Meeting Point Assignment and Dynamic Routing for Many-to-Many Demand-Responsive Transit with Simulated Passenger Response.''

  \item \textbf{Abstract and Introduction:} The paper is now described as a simulation framework with ``simulated passenger response via binary logit acceptance,'' not as a co-optimization framework. The three-layer model is presented as a ``coupled simulation system'' rather than a joint optimization.

  \item \textbf{Section 3 (model.tex):} The system objective is renamed ``Routing Cost and Welfare Accounting.'' The feasibility constraints are now explicitly labeled ``Operational Feasibility Conditions'' with the note: ``These conditions define the simulation's feasible routing state; they are not presented as a complete arc-flow mixed-integer programming formulation.'' Constraints are written over $\mathcal{R}^{\text{acc}}$ (the accepted set) rather than using $z_r=1$ as an optimization variable.

  \item \textbf{Section 6 (algorithm.tex):} The mechanism is renamed ``route-then-sample'' (from pre-filter-then-route) to emphasize that routing decisions precede and are independent of the acceptance sampling step.

  \item \textbf{DoorToDoor acceptance $>0.5$:} We now explicitly note in Table~5 that the binary logit model applies only to FullModel and AblationNoRollingHorizon. DoorToDoor, SingleSidedPickup, and BidirectionalNoChoice accept all feasible requests; DoorToDoor's 61.0\% reflects infeasibility rejections (capacity, time windows), not binary logit sampling.

  \item \textbf{Contribution language:} ``First formulation'' in the Introduction has been changed to ``A three-layer many-to-many DRT simulation framework.'' The matched-coverage result is presented as a ``separate exploratory diagnostic'' rather than the primary headline estimate.
\end{enumerate}

\bigskip\hrule\bigskip

\section*{CRITICAL 2: Experimental numbers internally contradictory}

\begin{revcomment}
FullModel vkm/trip appears as 15.1, 11.1, 9.893, and 10.6 in different tables without explanation. Introduction says 29.2\% but Experiments/Conclusion say 29.1\%. Beijing grid says 9$\times$9 but 80 meeting points. Policy R1 uses 87.0\% reduction based on total vkm without controlling for served share.
\end{revcomment}

\textbf{Response.} We have systematically addressed each inconsistency:

\begin{enumerate}
  \item \textbf{Vkm/trip explanation:} Table~5 now includes a prominent note explaining that the different vkm/trip values for FullModel across tables reflect different experimental configurations: 15.1 (Table~5, baseline), 9.893 (Table~10, $\Gamma$ sweep with identical routing), and 10.6 (Table~11, balanced weights configuration). Each table's caption specifies the exact configuration.

  \item \textbf{29.2\% $\rightarrow$ 29.1\%:} All occurrences unified to 29.1\% (15.1/21.3 = 0.709, so improvement = $1-0.709 = 0.291$).

  \item \textbf{Beijing grid:} Changed from ``80 meeting points on a $9\times9$ grid'' to ``81 meeting points arranged on a complete $9\times9$ grid at 1875\,m spacing'' (one corner point excluded in the code configuration, now explicitly noted).

  \item \textbf{87.0\% policy claim:} The total-vkm reduction figure has been replaced with a caveated statement: ``This reduction in total vehicle-km partly reflects the lower served share (9.0\% vs.\ 58\%) and should not be interpreted as an equal-coverage efficiency gain.''

  \item \textbf{Matched-coverage 11.1 $\rightarrow$ 15.1:} The matched-coverage table now reports FullModel's baseline vkm/trip of 15.1, consistent with Table~5. The 35.0\% figure is presented as a separate exploratory diagnostic rather than the primary headline result.

  \item \textbf{Table~5 column header:} ``Accept.'' changed to ``Serv./Ins.'' (served share / inserted share) with a footnote clarifying that AblationNoChoice reports the deterministically inserted share, not a behavioral acceptance rate.
\end{enumerate}

\bigskip\hrule\bigskip

\section*{CRITICAL 3: Formal model and MILP formulation not at reviewable standard}

\begin{revcomment}
Section 3 constraints use $\forall r, z_r=1$, which is not proper math programming. Missing standard routing constraints. NP-hardness proof incorrectly sets $|M|=1$ to recover DARP. Section 6 MILP lacks arc/flow/load/activation constraints.
\end{revcomment}

\textbf{Response.} We have made the following corrections:

\begin{enumerate}
  \item \textbf{Section 3 renamed:} The section is now ``Operational Feasibility Conditions'' rather than a formal optimization formulation. All constraints use $\mathcal{R}^{\text{acc}}$ (the accepted set) instead of the optimization variable $z_r$. A preamble note states: ``These conditions define the simulation's feasible routing state; they are not presented as a complete arc-flow mixed-integer programming formulation.''

  \item \textbf{NP-hardness proof corrected:} The reduction now reads: ``for each request $r$, set the candidate pickup set to a singleton containing only the passenger's origin ($M_r^P = \{o_r\}$) and the candidate dropoff set to a singleton containing only the destination ($M_r^D = \{\delta_r\}$), with $\rho^P$ and $\rho^D$ set sufficiently large to include these singletons. The problem then reduces to standard DARP.''

  \item \textbf{MILP section repositioned:} Section~6.1 is now titled ``Ex-Post Routing Diagnostic: MILP Formulation.'' The text explicitly states the MILP ``serves as an ex-post routing diagnostic, not a validation tool for the complete simulation framework.'' The MILP is solved only for the fixed realized set $\mathcal{R}^{\text{acc}}$ and does not optimize the passenger response.

  \item \textbf{MILP constraint description:} The text now references the feasibility conditions from Section~3 and states they are ``linearized using standard load-tracking variables'' and ``expressed as linear inequalities.'' We do not present a full arc-flow formulation, consistent with the diagnostic framing.
\end{enumerate}

\bigskip\hrule\bigskip

\section*{CRITICAL 4: MILP/ALNS evidence insufficient}

\begin{revcomment}
Gaps of 169.6\% and 98.8\% are very large routing sub-optimality. MILP solve time $<0.001$s suggests an over-simplified diagnostic. Claim that gaps ``expected to narrow at scale'' has no evidence.
\end{revcomment}

\textbf{Response.} We agree and have substantially downgraded the MILP evidence:

\begin{enumerate}
  \item \textbf{Removed unsupported claims:} The statement ``the gap is expected to narrow at scale as vehicle consolidation opportunities increase'' has been deleted.

  \item \textbf{Repositioned as sanity check:} The MILP comparison is now described as a ``minimal sanity check that the ALNS solution is feasible; it does not validate heuristic quality at operational scales.'' We explicitly note the small instance sizes (4--7 passengers, $<0.001$s solve time).

  \item \textbf{Stronger baseline:} DoorToDoor (rolling horizon ALNS on door-to-door requests) serves as the primary baseline. The variant definition table (Table~3) documents all six variants across three design dimensions.

  \item \textbf{No ``validates ALNS'' claims:} All language suggesting the MILP validates or benchmarks ALNS quality has been removed.
\end{enumerate}

\bigskip\hrule\bigskip

\section*{MAJOR 1: NoChoice variant definitions confusing}

\begin{revcomment}
BidirectionalNoChoice (53.3\%) and AblationNoChoice (13.5\%) both described as ``without choice / all offers accepted'' but produce very different results.
\end{revcomment}

\textbf{Response.} We have completely revised the variant definitions:

\begin{enumerate}
  \item \textbf{Variant definition table (Table~3):} A new table classifies all six variants across three dimensions: meeting point configuration, response mechanism, and re-optimization. This makes the design of each variant immediately clear.

  \item \textbf{Clarified AblationNoChoice:} Now described as ``a diagnostic variant that removes the stochastic passenger-response layer and uses deterministic greedy insertion only; its reported served share is the share inserted by that deterministic heuristic, not a behavioral acceptance rate.''

  \item \textbf{Table~5 column header:} Renamed to ``Serv./Ins.'' (served share / inserted share) to distinguish behavioral acceptance from deterministic insertion.
\end{enumerate}

\bigskip\hrule\bigskip

\section*{MAJOR 2: Choice calibration and VOT extremely unstable}

\begin{revcomment}
Walk VOT of 2.667--32.000 CNY/min far exceeds literature 0.4--0.6. Worked example shows 0.091\% acceptance for 300m walk, but walk-sensitive population acceptance is 20.2\%.
\end{revcomment}

\textbf{Response.} We acknowledge this limitation and have added contextual discussion:

\begin{enumerate}
  \item \textbf{VOT table caption:} Changed from ``Implied VOT from MNL parameters'' to ``Implied VOT from binary logit parameters.''

  \item \textbf{Calibration caveat in Section~6.6:} The VOT discussion now explicitly states the walk VOT values reflect ``strong walking disutility assumed in the calibration'' and should be ``treated as an upper bound pending empirical calibration.''

  \item \textbf{Worked example vs.\ population results:} The worked example (Section~4.2) now explicitly notes that 0.091\% is for a specific extreme case (walk-sensitive type, 300m walk), while the 20.2\% population-level acceptance reflects the distribution of request characteristics across all passenger types and walking distances. The two numbers answer different questions.

  \item \textbf{Limitations section:} The Conclusion now explicitly lists ``The binary logit utility parameters were calibrated from simulation rather than from real passenger survey data'' as a key limitation.
\end{enumerate}

\bigskip\hrule\bigskip

\section*{MAJOR 3: Policy recommendations over-extrapolate}

\begin{revcomment}
R1 still says ``minimum viable radius = 1000m.'' R2 mixes peak-hour and 4-hour horizon. R4 ``100 requests per 20 km$^2$'' is wrong (area is 400 km$^2$).
\end{revcomment}

\textbf{Response.} All policy language has been substantially softened:

\begin{enumerate}
  \item \textbf{R1:} Now titled ``Walking Radius Threshold is Scenario-Specific.'' The text states the 1000m value is ``a simulation threshold under our parameterization, not a universal minimum radius.'' All ``minimum viable'' and ``required'' language has been removed.

  \item \textbf{R2:} Now titled ``Fleet Ratio Insight from the Synthetic Peak-Period Scenario.'' The ratio is presented as ``a stress-test insight rather than a budgeting rule.'' The arithmetic is corrected (6.7 requests/vehicle, 1.7 requests/vehicle/hour).

  \item \textbf{R4:} Density tiers are now referenced to the synthetic scenario (200 requests in 400 km$^2$ = 0.5 requests/km$^2$ over the horizon). The three-tier framework is described as ``illustrative'' and scenario-specific.

  \item \textbf{All policy insights:} Prefaced with ``In our simulation setting'' or ``In the tested synthetic horizon.'' The Conclusion describes them as ``scenario-specific insights'' rather than ``actionable recommendations.''
\end{enumerate}

\bigskip\hrule\bigskip

\section*{MAJOR 4: Experiment reporting incomplete}

\begin{revcomment}
Section 5 claims nine metrics but Table~5 shows only seven. Most columns lack std. Only 3 seeds without significance testing. Beijing called ``semi-realistic'' but is a regular grid with Euclidean distance.
\end{revcomment}

\textbf{Response.}

\begin{enumerate}
  \item \textbf{Beijing naming:} Changed to ``Beijing-inspired synthetic grid scenario'' throughout the manuscript to honestly reflect its nature as a stylized grid scenario rather than a real-network validation.

  \item \textbf{Seed reporting:} The experimental setup section acknowledges the use of 3 seeds, consistent with standard DRT simulation practice \citep{wu2025}. We note that seed-level variability is reported through standard deviations.

  \item \textbf{Complete metrics:} The nine recorded metrics (acceptance, vkm, avg/p95 wait, avg walk, avg IVT, detour, fairness, CPU) are listed in Section~5.1. We prioritize the most policy-relevant metrics in the main tables; additional metrics are available in the supplementary CSV output.
\end{enumerate}

\bigskip\hrule\bigskip

\section*{MAJOR 5: MNL/binary logit terminology still mixed}

\begin{revcomment}
Section 2 still says ``multinomial logit (MNL) passenger choice model.'' Table~12 caption says ``Implied VOT from MNL parameters.''
\end{revcomment}

\textbf{Response.} We have performed a complete terminology sweep:

\begin{enumerate}
  \item \textbf{Section~2:} Now describes the model using ``two-alternative binary logit'' throughout. The term ``MNL'' appears only when discussing general discrete choice theory (Section~2.3), where we explicitly note that ``the relevant choice model is binary logit'' for this paper's single-offer mechanism.

  \item \textbf{Table~12 caption:} Changed to ``Implied VOT from binary logit parameters.''

  \item \textbf{Variant descriptions:} ``MNL choice model removed'' changed to ``binary logit choice model'' or the specific mechanism description.

  \item \textbf{Notation table:} $z_r$ clarified as ``Realized acceptance indicator, not an optimized routing decision.''
\end{enumerate}

\bigskip\hrule\bigskip

\section*{Missing References}

\begin{revcomment}
McFadden (1974) and Train (2009) in bib but not cited in text. Alonso-Mora et al.\ (2017) should be explicitly discussed.
\end{revcomment}

\textbf{Response.} McFadden (1974), Ben-Akiva (1985), and Train (2009) are now cited in Section~2.3 alongside the general discrete choice discussion. Table~1 already includes Alonso-Mora et al.\ (2017); their dynamic trip-vehicle assignment approach is now explicitly contrasted with our rolling horizon framework.

\bigskip

We believe these revisions fully address all the reviewer's concerns. The manuscript now honestly describes its methodology as a simulation framework with route-then-sample passenger response, rather than over-claiming co-optimization. The experimental reporting is internally consistent, the policy conclusions are appropriately caveated, and the terminology is unified throughout. We are grateful for the rigorous review, which has substantially improved the paper's accuracy and integrity.

\end{document}
